% Generated by `python -m experiments.experiment_2.appendix_captions` -- do not edit by hand.
% Methods/appendix text for fig:exp2-main (main_figure.tex). Place it under a heading
% carrying the label, e.g. \subsection{Posterior predictive checks: details}\label{app:exp2-main-details}
% TODO(verify): RDM dataset citation; Kvam (2019) for the CDM data and response apparatus.
The posterior predictive in Figure~\ref{fig:exp2-main} is built from 500 posterior draws per subject. Each draw simulates that subject's own trials under their own design (the contamination process included), and the simulated trials are pooled over subjects, giving one predictive data set per draw; quantiles and choice proportions are computed per data set, and the band is their central 95\% interval across draws. Data: DDM, Stump et al. (2026), 72 subjects with at least 100 trials (100--116, median 112 trials per subject), regressors stimulus type and factual truth (effect-coded $\pm 0.5$) and their interaction; RDM, a three-alternative Stroop task, 149 subjects (171--192, median 189 trials), congruency and session ($\pm 0.5$) and their interaction, with stimulus-coded responses (the ink colour is the favoured accumulator); CDM, 12 subjects (456--599, median 476 trials), speed instruction ($\pm 0.5$) and three difficulty levels (dummy-coded against the easiest) with their interaction. The CDM's recorded response deviation spans $[-\pi/2, \pi/2]$ and is doubled to the model's response angle on $[-\pi, \pi]$, i.e.\ treated as a half-angle. 
